% 2.1.2
% Ficheiro incluído por main.tex

\subsection[Ubuntu: IP estático e DNS do Windows]{Ubuntu: IP estático e DNS do Windows}

\passo{Passo 1 - Iniciar a VM e abrir o terminal}

Trabalhar \textbf{na consola da VM} (não por SSH).

\passo{Passo 2 - Ver a configuração de rede existente}

\begin{lstlisting}[style=bashstyle]
ls /etc/netplan
nmcli -f NAME,DEVICE,TYPE connection show
\end{lstlisting}

\begin{itemize}
\item \texttt{ls /\allowbreak{}etc/\allowbreak{}netplan} lista os ficheiros do Netplan.
\item \texttt{nmcli ... connection show} lista as ligações geridas pelo NetworkManager.
\end{itemize}

\passo{Passo 3 - Criar uma pasta de backup}

\begin{lstlisting}[style=bashstyle]
sudo mkdir -p /root/netplan-bak
\end{lstlisting}

\passo{Passo 4 - Remover os ficheiros em conflito}

\begin{lstlisting}[style=bashstyle]
cd /etc/netplan
sudo mv 00-installer-config.yaml 01-netcfg.yaml 90-NM-*.yaml /root/netplan-bak/
ls /etc/netplan
cd ~
\end{lstlisting}

\emph{Verificação esperada:} o \texttt{ls} mostra \textbf{só} \texttt{01-network-manager-all.yaml}.

\passo{Passo 5 - Aplicar a limpeza}

\begin{lstlisting}[style=bashstyle]
sudo netplan apply
\end{lstlisting}

\passo{Passo 6 - Garantir que o NetworkManager está a correr}

\begin{lstlisting}[style=bashstyle]
sudo systemctl start NetworkManager
systemctl is-active NetworkManager
\end{lstlisting}

\emph{Verificação esperada:} \texttt{active}.

\passo{Passo 7 - Criar a ligação com IP estático e DNS do Windows Server}

\begin{lstlisting}[style=bashstyle]
sudo nmcli connection add type ethernet ifname ens18 con-name labX \
     ipv4.method manual \
     ipv4.addresses <IP_CLIENTE>/24 \
     ipv4.gateway 192.168.X.1 \
     ipv4.dns <IP_DC> \
     ipv4.dns-search isec.local \
     ipv4.ignore-auto-dns yes
\end{lstlisting}

\emph{Verificação esperada:} o comando termina sem erros.

\passo{Passo 8 - Ativar a ligação e confirmar}

\begin{lstlisting}[style=bashstyle]
sudo nmcli connection up labX
nmcli -f NAME,DEVICE,TYPE connection show
\end{lstlisting}

\emph{Verificação esperada:} o perfil \texttt{labX} aparece associado à placa \texttt{ens18}.

\begin{lstlisting}[style=bashstyle]
sudo nmcli connection delete 'NOME'
\end{lstlisting}

\passo{Passo 9 - Confirmar o IP e a rota}

\begin{lstlisting}[style=bashstyle]
ip -br a
ip route
\end{lstlisting}

\emph{Verificação esperada:}

\begin{itemize}
\item \texttt{ens18} aparece \textbf{UP} com \texttt{\textless{}IP\_CLIENTE\textgreater{}/\allowbreak{}24} - e \textbf{não} com um endereço \texttt{169.254.x.x};
\item existe \textbf{uma única} linha \texttt{default via 192.168.X.1} (o conflito do início desapareceu).
\end{itemize}

\passo{Passo 10 - Confirmar que o DNS é o do Windows}

\begin{lstlisting}[style=bashstyle]
resolvectl status ens18
\end{lstlisting}

\emph{Verificação esperada:} \texttt{Current DNS Server: \textless{}IP\_DC\textgreater{}}, \texttt{DNS Servers: \textless{}IP\_DC\textgreater{}} (sem \texttt{8.8.8.8} nem \texttt{1.1.1.1}), \texttt{DNS Domain: isec.local}.

\passo{Passo 11 - Testar a ligação ao servidor}

\begin{lstlisting}[style=bashstyle]
ping -c 3 <IP_DC>
\end{lstlisting}

\emph{Verificação esperada:} \texttt{0\% packet loss}.

\passo{Passo 12 - Testar o DNS do domínio}

\begin{lstlisting}[style=bashstyle]
dig @<IP_DC> isec.local +short
dig @<IP_DC> _ldap._tcp.dc._msdcs.isec.local SRV +short
\end{lstlisting}

\emph{Verificação esperada:}

\begin{itemize}
\item o 1.º devolve \texttt{\textless{}IP\_DC\textgreater{}};
\item o 2.º devolve algo como \texttt{0 100 389 sr2-srv.isec.local.}
\end{itemize}

\passo{Passo 13 - Corrigir o ping a \texttt{isec.local} (desativar o mDNS para \texttt{.local})}

\begin{lstlisting}[style=bashstyle]
sudo cp /etc/nsswitch.conf /etc/nsswitch.conf.bak
grep ^hosts /etc/nsswitch.conf
sudo sed -i 's/mdns4_minimal \[NOTFOUND=return\] //' /etc/nsswitch.conf
grep ^hosts /etc/nsswitch.conf
\end{lstlisting}

\emph{Verificação esperada:} a linha passa de \texttt{hosts: files mdns4\_minimal [NOTFOUND=return] dns} para \texttt{hosts: files dns}.

\passo{Passo 14 - Testar o \texttt{isec.local}}

\begin{lstlisting}[style=bashstyle]
getent hosts isec.local
ping -c 3 isec.local
\end{lstlisting}

\emph{Verificação esperada:} \texttt{getent} devolve \texttt{\textless{}IP\_DC\textgreater{} isec.local} e o ping responde com \texttt{0\% packet loss}.

